\BourbakiExercises{习题}

\begin{enumerate}
\def\labelenumi{\arabic{enumi})}
\tightlist
\item
  设 \(\mathrm{K}\) 是一个交换域。设 \(\mathrm{A}\) 是一个 K-代数，它的基由单位元与两个元素 \(e, f\) 组成，满足 \(e^2 = e\)，\(ef = f\)，\(fe = f^2 = 0\)。a) 证明 \(\mathrm{A}\) 的正则表示与余正则表示不同构（计算 \(\operatorname{Tr}_{\mathrm{A / K}}(e)\) 与 \(\operatorname{Tr}_{\mathrm{A}^{\circ} / \mathrm{K}}(e)\)）。
\end{enumerate}

\begin{enumerate}
\def\labelenumi{\alph{enumi})}
\setcounter{enumi}{1}
\tightlist
\item
  证明正则表示包含一个子表示，它既不同构于余正则表示的任何子表示，也不同构于余正则表示的任何商表示。此外，存在 A 的一个单表示，它不同构于正则表示的任何子表示。
\end{enumerate}

\begin{enumerate}
\def\labelenumi{\arabic{enumi})}
\setcounter{enumi}{1}
\tightlist
\item
  设 \(A\) 是一个 K-代数。我们把对偶 \(A^*\)（\(A\) 的对偶）视为一个左 \(A\)-模。
\end{enumerate}

\begin{enumerate}
\def\labelenumi{\alph{enumi})}
\item
  映射 \(E \mapsto E'\) 是从 \(A^{*}\) 的 A-子模所成的集合到 A 的右理想所成的集合上的一个满射。当 E 是这样一个子模时，\(E'\) 在 \(A^{*}\) 中的正交是 \(A^{*}\) 的一个包含 E 的子模。
\item
  证明 \(\mathbf{A}^*\) 的底座在 A 中的正交是 A 的根。
\item
  证明在 \(A^{*}\) 中，\(A_{d}\) 的底座的正交是 A-模 \(A^{*}\) 的根。由此推出：\(A^{*}\) 是半单 A-模，当且仅当环 A 是半单的。
\end{enumerate}

¶ 3) 设 A 是一个阿廷环，\(\overline{\mathrm{A}}\) 是环 \(\mathrm{A} / \Re (\mathrm{A})\)。我们取 A 的一个分解，把它分解为不可分解理想 \(\mathrm{Ae}_i\)（\(i\in \mathrm{I}\)）的直和，其中诸 \(e_i\) 是正交幂等元（VIII, 第 179 页，习题 6, d)）。用 \(\mathrm{I} = \cup_{\alpha \in \mathrm{R}}\mathrm{I}_{\alpha}\) 表示 I 的这样一个分拆：I 的两个元素 \(i,j\) 属于同一个集合 \(\mathrm{I}_{\alpha}\)，当且仅当模 \(\mathrm{Ae}_i\) 与 \(\mathrm{Ae}_j\) 同构。\(\overline{\mathrm{A}}\)-模 \(\bigoplus_{i\in \mathrm{I}_{\alpha}}\overline{\mathrm{A}}\overline{e}_i\) 是 \(\overline{\mathrm{A}}\) 的同型分量（VIII, 第 179 页，习题 6, d) 与 7, a)）

\begin{enumerate}
\def\labelenumi{\alph{enumi})}
\item
  证明：若环 A 是自内射的（VIII, 第 53 页，习题 8），则存在 R 的一个置换 \(\sigma\)，使得每个右理想 \(e_i\mathrm{A}\)（对 \(i\in \mathrm{I}_{\alpha}\)）都包含唯一一个同构于 \(\overline{e}_j\overline{\mathrm{A}}\) 的极小右理想（对 \(j\in \mathrm{I}_{\sigma (\alpha)}\)）（注意幂等元 \(e\) 的右零化理想是 \((1 - e)\mathrm{A}\)）。由此推出 A 的左底座 \(\mathfrak{s}\) 包含于其右底座 \(\mathfrak{t}\)（证明每个右理想 \(e_i\mathfrak{s}\) 都是极小的，为此证明它是一个极大左理想的零化理想，参见 VIII, 第 178 页，习题 2）。最后断定 \(\mathfrak{s} = \mathfrak{t}\)，并且每个左理想 \(\mathrm{Ae}_i\)（对 \(i\in \mathrm{I}_{\sigma (\alpha)}\)）都包含唯一一个同构于 \(\overline{\mathrm{A}}\overline{e}_j\) 的极小左理想（对每个 \(j\in \mathrm{I}_{\alpha}\)）。
\item
  反之，若每个左理想 \(Ae_{i}\) 都包含唯一的极小左理想，且每个右理想 \(e_{i}A\) 都包含唯一的极小右理想，则环 A 是自内射的。（利用 VIII, 第 179 页的习题 7 证明 s 的每个足都是一个极小双边理想。由此推出 t 包含 s，进而 s = t。最后证明 VIII, 第 53 页习题 8 的条件 \(\beta\) ) 得到满足。）
\item
  设 A 是一个交换域上的有限次数代数，并且是一个自内射环。设 \(\alpha \in R\)，\(i \in I_{\alpha}\)，\(j \in I_{\sigma(\alpha)}\)。证明与 A-模 \(Ae_{j}\) 相伴的 A 的线性表示同构于 \(A^{\circ}\) 的线性表示（与右 A-模 \(e_{i}A\) 相伴的线性表示）的转置。（设 \(U_{i}\) 是 \(A^{*}\) 的作为 \((1 - e_{i})A\) 的正交的那个左子模。由 VIII, 第 179 页的习题 7, d) 推出 \(U_{i}\) 同构于 \(Ae_{j}\) 的一个商模，从而有 long \(e_{i}A \leqslant long e_{j}A\)。断定这两个长度相等，于是 \(U_{i}\) 同构于 \(Ae_{j}\)。利用 b) 与习题 2, a) 证明该命题的一个逆命题。）
\end{enumerate}

\begin{enumerate}
\def\labelenumi{\arabic{enumi})}
\setcounter{enumi}{3}
\tightlist
\item
  保留习题 3 的假设。
\end{enumerate}

\begin{enumerate}
\def\labelenumi{\alph{enumi})}
\item
  证明：若 A 是一个弗罗贝尼乌斯环（VIII, 第 53 页，习题 9），则有 \(\operatorname{Card}(\mathrm{I}_{\alpha}) = \operatorname{Card}(\mathrm{I}_{\sigma(\alpha)})\)（对每个 \(\alpha \in R\)）。并证明其逆。
\item
  一个交换域上的有限次数代数是弗罗贝尼乌斯环，当且仅当它的正则表示同构于它的余正则表示（利用习题 3, c)）。
\end{enumerate}

\begin{enumerate}
\def\labelenumi{\arabic{enumi})}
\setcounter{enumi}{4}
\tightlist
\item
  设 K 是一个交换域，A 是一个 K-代数。我们赋予 \(A^{*}\) 其典范左 A-模结构。对 \(t \in A^{*}\)，用 \(\varphi_{t}\) 表示映射 \((x, y) \mapsto t(xy)\)（从 \(A \times A\) 到 K 的映射）。
\end{enumerate}

\begin{enumerate}
\def\labelenumi{\alph{enumi})}
\item
  证明这定义了一个 K-线性同构，从 \(A^{*}\) 到由 K-双线性映射 \(\varphi: A \times A \to K\)（满足 \(\varphi(xy, z) = \varphi(x, yz)\)，对 \(x, y, z \in A\)）所成的集合。
\item
  我们说元素 \(t \in A^{*}\) 是可逆的，如果同态 \(a \mapsto at\)（从 A 到 \(A^{*}\)）是双射。设 A 是一个有限次数的 K-代数。证明下列条件等价：
\end{enumerate}

\begin{enumerate}
\def\labelenumi{(\roman{enumi})}
\item
  A 是弗罗贝尼乌斯环（习题 4）。
\item
  A 容许一个可逆线性型。
\item
  A 容许一个核不包含 \(\mathbf{A}\) 的任何左理想的线性型。
\end{enumerate}

\begin{enumerate}
\def\labelenumi{\arabic{enumi})}
\setcounter{enumi}{5}
\tightlist
\item
  设 K 是一个交换域。我们说 K-代数 A 是对称的，如果它容许一个可逆的迹型线性型（VIII, 第 115 页，习题 3）。
\end{enumerate}

\begin{enumerate}
\def\labelenumi{\alph{enumi})}
\item
  证明对称代数是有限次数的且是弗罗贝尼乌斯环。
\item
  证明 \(\mathrm{K}\) 上有限群的群代数是对称的。
\item
  \(\mathrm{K}\) 上有限次数的半单代数是对称的。
\end{enumerate}

